\section{Introduction}
\label{sec:intro}

Probabilistic Graphical Models (PGMs) provide a language for describing the conditional (in)dependence structure of complex systems, but there is a fundamental tension between \textbf{generative (causal)} and \textbf{discriminative (inference) factorization}: a single model cannot serve both mechanistic simulation and posterior inference efficiently---for structured flows we make this precise in Theorem~\ref{thm:marginalization}. A causal model $(X \to Y)$ excels at interventional queries $p(Y | do(X))$ but requires expensive inversion or Markov chain Monte Carlo (MCMC) for inference $p(X | Y)$; an inverse model $(Y \to X)$ enables instant amortized inference but breaks the causal link, rendering it blind to the mechanisms of the world.

A clinical model of a disease system must support both treatment planning---``what will happen to this biomarker if we administer the drug?'', an interventional query posed \emph{along} the causal direction---and diagnosis---``which disease state best explains these symptoms?'', a posterior query posed \emph{against} it. The same pairing recurs across the sciences (a materials simulator predicts microstructure from processing parameters; the experimentalist needs the parameters that explain an observed microstructure).

The root of this tension is twofold. First, by the \textbf{Principle of Independent Mechanisms (PIM)}, the mechanism generating an effect, $P(\text{Effect}|\text{Cause})$, is independent of the distribution of the cause \citep{Peters2017}: a flow optimized for the causal direction learns a modular, reusable mechanism, but inverting it for inference entangles these components via Bayes' rule, so forcing one set of parameters to capture both directions compromises both. Second, causal models face a \textbf{structural mismatch} when used for inference: causal directed acyclic graphs (DAGs) are often sparse, containing ``colliders'' ($X_1 \to Y \leftarrow X_2$) whose parents are marginally independent, yet observing the child $Y$ induces a dependency between them (``explaining away'' \citep{Koller2009}). The posterior is strictly denser than the prior, and a flow built with the sparsity of the causal DAG lacks the edges required to represent these posterior correlations.

We propose \textbf{StructFlow} (\textbf{Struct}ural \textbf{Flow} Alignment), a framework for Continuous Normalising Flows that provides a unified \emph{interface} and \emph{methodology} for causal reasoning and statistical inference. By formalizing the mapping between PGM structures and flow architectures, we allow a user to define the system's dependencies once and ``compile'' two distinct model artifacts depending on their needs: a \textbf{World Model} (Causal Alignment) for mechanistic simulation, and an \textbf{Inference Engine} (Inference Alignment) that employs moralization to capture posterior correlations. Importantly, we train \emph{separate} models for each task. Our goal is to demonstrate these algorithmic ideas in principle---one specification, compiled automatically for arbitrary DAGs, answering causal and inference queries correctly---so we evaluate on problems with exact ground truth, and to characterize where the design choices (structure, conditioning mechanism, backbone) genuinely matter.

Our contributions are four-fold:
\begin{enumerate}
    \item A \textbf{Unified Methodology} that automates the compilation of arbitrary PGMs into structured continuous normalising flow (CNF) architectures, supporting both causal (generative) and inference (discriminative) tasks through a common interface (Table~\ref{tab:capability}).
    \item The \textbf{Implementation} of both Static and Dynamic conditioning for arbitrary DAGs---generalizing block-triangular flows from a single observation-to-target split to arbitrary graphical structure---with an \textbf{analysis of dynamic conditioning's scaling}: when and why its parallel sampling degrades (deep dependency chains with informative evidence; Section~\ref{sec:chain-stiffness}).
    \item An \textbf{Empirical Characterization} of why structure matters: \rev{the posterior dependencies added by the compiler are necessary to represent collider posteriors} (Section~\ref{sec:moralization-ablation}); the compiled inference engine amortizes, answering posterior queries roughly $11\times$ faster than per-query MCMC beyond a small crossover (Section~\ref{sec:amortization}); and these capabilities follow from the compiled structure, not the choice of continuous-time generative model (Section~\ref{sec:backbone-ablation}).
    \item A \textbf{Recovery Analysis} delimiting what the compiled causal model provably identifies: the causally-aligned flow \emph{is} a structural causal model on the given graph (Lemma~\ref{lem:mech-main}), so fitting the observational joint recovers every interventional distribution for arbitrary mechanisms (Theorem~\ref{thm:level2}), and counterfactuals exactly for mechanisms invertible in a scalar noise or additive in a multivariate one (Theorem~\ref{thm:level3}). What this does \emph{not} give is the graph itself: it is unidentifiable in the linear--Gaussian case (Proposition~\ref{prop:mec}), so structure discovery is the single assumption separating the framework from full causal identification.
\end{enumerate}